\chapter{Methodology}
\label{ch:methodology}

\section{Introduction}
\label{sec:meth_intro}
This chapter delineates the technical methodology, architectural blueprints, database models, and operational workflows implemented to realize the web-based government task tracking and alert platform. The system is engineered using a decoupled, multi-tier client-server paradigm emphasizing scalability, modularity, and strict cryptographic security.

\section{Overall System Architecture}
\label{sec:system_architecture}
The system architecture comprises three primary tiers:
\begin{enumerate}
    \item \textbf{Presentation Layer (Client):} Built using React.js for desktop web browsers and React Expo for cross-platform mobile accessibility.
    \item \textbf{Application Layer (Server):} Powered by a stateless Node.js / Express.js REST API providing endpoint routing, authentication middleware, task business logic, and scheduler services.
    \item \textbf{Data Persistence Layer:} Implemented via a normalized MySQL relational database engine configured with ACID-compliant transactions and foreign key referential integrity.
\end{enumerate}

\section{Role-Based System Design}
\label{sec:rbac_design}
To match the organizational hierarchy of public administration, four distinct functional roles were designed:

\subsection{Admin Module}
The Administrator possesses supervisory authority over platform metadata. Capabilities include:
\begin{itemize}
    \item Provisioning and deactivating user accounts across departments.
    \item Configuring departmental schemas, designation hierarchies, and statutory holiday calendars.
    \item Inspecting global system security logs, database health metrics, and audit trails.
\end{itemize}

\subsection{Nodal Officer Module}
Nodal Officers represent central directing authorities (e.g., District Magistrates or Department Directors). They are empowered to:
\begin{itemize}
    \item Author new statutory directives and assign cross-departmental tasks.
    \item Define priority tiers (Low, Medium, Critical) and statutory completion deadlines.
    \item Inspect executive aggregations of departmental progress and initiate direct escalations.
\end{itemize}

\subsection{Department Module}
Department Heads receive broad administrative directives and decompose them into actionable subtasks. Features include:
\begin{itemize}
    \item Delegating departmental tasks to designated executive officers.
    \item Reviewing completion proofs and intermediate deliverables submitted by officers.
    \item Requesting statutory deadline extensions or reassignments from the Nodal authority.
\end{itemize}

\subsection{Officer Module}
Field and executive officers execute operational duties. Their module provides:
\begin{itemize}
    \item An interactive personal task backlog sorted by urgency and deadline.
    \item File upload interfaces for digital proof of work (PDFs, geo-tagged imagery).
    \item Direct status toggle mechanisms (\textit{Pending}, \textit{In Progress}, \textit{Submitted for Verification}).
\end{itemize}

\section{Frontend Design Using React JS and React Expo}
\label{sec:frontend_design}
The client interface is engineered using React.js (v18+) with modern functional component architecture, React Hooks, and Tailwind CSS. React Router DOM manages client-side protected route boundaries. State synchronization utilizes React Query for server cache invalidation and Axios for HTTP communications. For mobile accessibility by field staff, React Expo enables shared TypeScript interface models across Android and iOS deployments.

\section{Backend Implementation Using Node.js and Express}
\label{sec:backend_implementation}
The backend application runtime leverages Node.js (v20 LTS). Express.js serves as the middleware framework. The service layer enforces architectural separation of concerns into:
\begin{itemize}
    \item \textbf{Routers:} Route definition and parameter sanitization.
    \item \textbf{Controllers:} Orchestration of request payload parsing and response generation.
    \item \textbf{Services:} Pure business logic execution and workflow calculations.
    \item \textbf{Data Access Layer (DAO):} Parameterized SQL query execution preventing SQL-injection vulnerabilities.
\end{itemize}

\section{Database Design Using MySQL}
\label{sec:database_design}
The relational schema is normalized to Third Normal Form (3NF) to eliminate data redundancy and preserve integrity. Key entities include \texttt{Users}, \texttt{Roles}, \texttt{Departments}, \texttt{Tasks}, \texttt{Task\_Assignments}, \texttt{Submissions}, and \texttt{Notifications}.

\begin{table}[ht]
\centering
\caption{Core Database Schema Entities and Relational Mappings}
\label{tab:db_schema}
\begin{tabularx}{\textwidth}{l l X}
\toprule
\textbf{Table Name} & \textbf{Primary Key} & \textbf{Key Foreign References \& Attributes} \\
\midrule
\texttt{users} & \texttt{user\_id} (INT) & \texttt{department\_id}, \texttt{role\_id}, email, password\_hash \\
\texttt{roles} & \texttt{role\_id} (INT) & role\_name, permission\_flags \\
\texttt{departments} & \texttt{dept\_id} (INT) & dept\_name, nodal\_officer\_id \\
\texttt{tasks} & \texttt{task\_id} (INT) & creator\_id, priority, deadline, status \\
\texttt{task\_logs} & \texttt{log\_id} (INT) & \texttt{task\_id}, actor\_id, previous\_status, new\_status, timestamp \\
\bottomrule
\end{tabularx}
\end{table}

\section{Role-Based Authentication and Authorization}
\label{sec:auth_auth}
Authentication is governed by JSON Web Tokens (JWT). Upon successful credential verification using bcrypt (salt rounds = 12), the server generates a cryptographically signed HMAC-SHA256 token containing the user identity, department ID, and role claim. Custom Express middleware verifies this token on protected routes and validates role access privileges prior to controller invocation.

\section{Alert and Task Assignment Workflow}
\label{sec:workflow}
An automated cron scheduler executes periodically to inspect active task deadlines. The alert engine classifies deadlines into three temporal windows:
\begin{itemize}
    \item \textbf{Window 1 ($t \le 48\text{ hrs}$):} Standard reminder dispatched to assignee.
    \item \textbf{Window 2 ($t \le 24\text{ hrs}$):} High-priority reminder dispatched to assignee and Department Head.
    \item \textbf{Window 3 ($t < 0\text{ hrs}$):} Overdue escalation triggered; notification dispatched to Nodal Officer with automated incident tagging.
\end{itemize}

\section{Report Submission and Status Tracking Mechanism}
\label{sec:report_submission}
Officers submit task completion documentation via multipart/form-data endpoints. Uploaded assets undergo MIME-type validation and are stored in an encrypted file repository. A state machine governs transitions:
\begin{equation}
    \text{ASSIGNED} \longrightarrow \text{IN\_PROGRESS} \longrightarrow \text{SUBMITTED} \longrightarrow \text{VERIFIED} \mid \text{REJECTED}
\end{equation}

\section{Summary of Methodology}
\label{sec:meth_summary}
The adopted methodology establishes an auditable, responsive, and robust software architecture. By decoupling presentation from business services and grounding data transactions in an ACID-compliant MySQL database, the platform satisfies governmental security, transparency, and operational velocity requirements.
